\section{Methodology}
\label{sec:method}

\subsection{Problem Setup}

We consider a training dataset $\mathcal{D} = \{(x_i, y_i, g_i)\}_{i=1}^N$ where group labels $g_i$ are hidden during training and evaluation. ERM minimizes $\mathcal{L}(\theta) = \frac{1}{N}\sum_i \ell(f_\theta(x_i), y_i)$ without observing $g_i$. We ask: at training time, does any per-sample statistic derived from the ERM loss surface discriminate minority from majority group membership --- using only $f_\theta(x_i)$ and $\ell$, without observing $g_i$?

\subsection{Per-Sample Last-FC Hessian Trace}

For sample $i$ at training checkpoint $t$, we define the per-sample last-fc Hessian trace as:
\begin{equation}
\text{Tr}(H_i^{fc}) = \text{Tr}\!\left(\frac{\partial^2 \ell(f_{\theta_t}(x_i), y_i)}{\partial W_{fc}^2}\right)
\end{equation}
where $W_{fc} \in \mathbb{R}^{2048\times 2}$ is the weight matrix of the last fully-connected layer of ResNet-50. We restrict to the last-fc layer for three reasons: (1) tractability --- 4098 parameters vs.\ 25M total; (2) architectural motivation --- the classification head directly encodes the spurious feature alignment learned by the backbone; (3) mechanistic interpretability --- the per-sample last-fc Hessian has the closed-form decomposition $\text{Tr}(H_i^{fc}) \propto \|x_i\|^2 p_i(1-p_i)$ for cross-entropy loss, isolating the feature-norm and confidence contributions.

\subsection{Hutchinson Trace Estimation}

Computing $\text{Tr}(H)$ exactly requires full Hessian materialization. We use Hutchinson's estimator \cite{avron2011randomized}:
\begin{equation}
\text{Tr}(H) \approx \frac{1}{K} \sum_{k=1}^K v_k^\top H v_k
\end{equation}
where $v_k \in \{{\pm}1\}^d$ are Rademacher random vectors, and $Hv_k = \frac{\partial}{\partial W}\!\left(\frac{\partial \ell}{\partial W} \cdot v_k\right)$ is computed as a second-order reverse-over-reverse automatic differentiation pass. We implement this via \texttt{torch.func.vmap} + \texttt{torch.func.vjp} on a per-sample basis, enabling vectorized computation across samples without storing the full Hessian or Jacobian.

\textbf{Choice of $K{=}50$:} Prior diagnostic work \cite{yao2020pyhessian} uses $K{=}256$ for batch-level Hessian traces on large networks. For the last-fc layer ($d{=}4098$), $K{=}50$ Rademacher probes provide trace estimates with coefficient of variation CV$<$3\% across all seeds in our experiments --- confirmed as the minimum stable setting for this architecture.

\subsection{Training Protocol}

We train ResNet-50 (torchvision, ImageNet pretrained) on Waterbirds v1.0 with ERM:
\begin{itemize}
\item Optimizer: SGD (lr$=3{\times}10^{-3}$, momentum$=0.9$, weight decay$=1{\times}10^{-4}$)
\item Epochs: 50
\item Checkpoints: $t \in \{0, 1, 5, 10, 20, 50\}$
\item Seeds: 5 independent seeds (seeds 1--5)
\end{itemize}

The $t{=}0$ checkpoint is the pretrained ImageNet initialization \emph{before any Waterbirds training}. This serves as the epoch-0 control for distinguishing ERM-induced signals from pretrained artifacts (see Section~\ref{sec:epoch0}).

\textbf{Dataset:} Waterbirds v1.0 contains 4795 training samples with four groups defined by (bird type, background): (waterbird/water, waterbird/land, landbird/water, landbird/land). Groups 1+3 (minority groups) comprise 240 samples (5\% of training set). Group labels are used only for post-hoc evaluation; they are not observed during training.

\subsection{Trace Ratio and $t^*$ Selection}

At each checkpoint $t$, we compute the trace ratio:
\begin{equation}
R(t) = \frac{\overline{\text{Tr}(H^{fc})}_{i \in \text{minority}}}{\overline{\text{Tr}(H^{fc})}_{i \in \text{majority}}}
\end{equation}

The optimal checkpoint $t^* = \arg\max_t R(t)$ selects the epoch of maximum minority-majority trace divergence. This is the evaluation checkpoint for all discriminability metrics. $R(t)$ is computed without observing group labels at training time --- only the rank ordering of traces is used for AUROC evaluation.

\subsection{Evaluation Protocol}
\label{sec:epoch0}

\textbf{Existence gate (H-E3):} At checkpoint $t^*$, we compute AUROC for predicting minority membership ($g_i \in \{1,3\}$) from the scalar trace $\text{Tr}(H_i^{fc}(t^*))$. Three simultaneous criteria must hold:
\begin{enumerate}
\item $\text{AUROC}(t^*) \geq 0.85$
\item $\text{AUROC}(t{=}0) < 0.70$ (ERM-induction control)
\item Spearman $\rho(t, \text{AUROC}(t)) \geq 0.8$ on the rising AUROC segment
\end{enumerate}

The epoch-0 control (criterion 2) is mandatory: it rules out that minority-majority trace differences are inherited from ImageNet pretraining rather than induced by ERM training on Waterbirds.

\textbf{Mechanism test (H-M1):} At $t^*$, we measure the mean training-set prediction confidence for minority and majority samples separately. The confidence-differential mechanism predicts $p_{\text{minority}}(t^*) \in [0.3, 0.7]$ (boundary-condition proximity) in $\geq$3/5 seeds. This is the pre-registered gate for mechanism verification.

\textbf{Reproducibility standard:} All 5 seeds are evaluated independently. The H-E3 gate requires $\geq$4/5 seeds to pass all three criteria simultaneously.
